% Options for packages loaded elsewhere
\PassOptionsToPackage{unicode}{hyperref}
\PassOptionsToPackage{hyphens}{url}
\documentclass[
]{article}
\usepackage{xcolor}
\usepackage{amsmath,amssymb}
\setcounter{secnumdepth}{-\maxdimen} % remove section numbering
\usepackage{iftex}
\ifPDFTeX
  \usepackage[T1]{fontenc}
  \usepackage[utf8]{inputenc}
  \usepackage{textcomp} % provide euro and other symbols
\else % if luatex or xetex
  \usepackage{unicode-math} % this also loads fontspec
  \defaultfontfeatures{Scale=MatchLowercase}
  \defaultfontfeatures[\rmfamily]{Ligatures=TeX,Scale=1}
\fi
\usepackage{lmodern}
\ifPDFTeX\else
  % xetex/luatex font selection
\fi
% Use upquote if available, for straight quotes in verbatim environments
\IfFileExists{upquote.sty}{\usepackage{upquote}}{}
\IfFileExists{microtype.sty}{% use microtype if available
  \usepackage[]{microtype}
  \UseMicrotypeSet[protrusion]{basicmath} % disable protrusion for tt fonts
}{}
\makeatletter
\@ifundefined{KOMAClassName}{% if non-KOMA class
  \IfFileExists{parskip.sty}{%
    \usepackage{parskip}
  }{% else
    \setlength{\parindent}{0pt}
    \setlength{\parskip}{6pt plus 2pt minus 1pt}}
}{% if KOMA class
  \KOMAoptions{parskip=half}}
\makeatother
\setlength{\emergencystretch}{3em} % prevent overfull lines
\providecommand{\tightlist}{%
  \setlength{\itemsep}{0pt}\setlength{\parskip}{0pt}}
\usepackage{bookmark}
\IfFileExists{xurl.sty}{\usepackage{xurl}}{} % add URL line breaks if available
\urlstyle{same}
\hypersetup{
  hidelinks,
  pdfcreator={LaTeX via pandoc}}

\author{}
\date{}

\begin{document}

\subsubsection{\texorpdfstring{\textbf{1. Project
Overview}}{1. Project Overview}}\label{project-overview}

\textbf{Problem Statement:}\\
Financial institutions face a growing backlog of suspicious transactions
flagged by rule-based systems. Analysts manually review these alerts by
collecting data from multiple sources (e.g., online portals, mainframe
systems), a process that is \textbf{time-consuming, error-prone, and
inefficient}. With transaction volumes rising, the manual review process
cannot scale, leading to delays in detecting genuine fraud.

\textbf{Need for Automation:}\\
The goal is to \textbf{automate data collection} and \textbf{prioritize
high-risk transactions} using machine learning (ML), reducing analyst
workload and accelerating decision-making.

\begin{center}\rule{0.5\linewidth}{0.5pt}\end{center}

\subsubsection{\texorpdfstring{\textbf{2. Process
Analysis}}{2. Process Analysis}}\label{process-analysis}

\textbf{Existing Workflow:}\\
1. An analyst receives a flagged transaction.\\
2. \textbf{Manual Data Collection:}\\
- Log into 3-4 web portals (e.g., customer profiles, transaction
history).\\
- Query mainframe systems via terminal emulators (e.g., AS400).\\
- Copy-paste data into Word/Excel for review.\\
3. \textbf{Time Breakdown:}\\
- Portal navigation: \textbf{8--10 minutes/transaction}.\\
- Data validation: \textbf{5--7 minutes/transaction}.\\
- Total: \textbf{15--20 minutes/transaction}.

\textbf{Inefficiencies:}\\
- Repetitive tasks dominate 80\% of the workflow.\\
- Analysts handle only \textbf{30--40 alerts/day}, leading to backlogs.

\begin{center}\rule{0.5\linewidth}{0.5pt}\end{center}

\subsubsection{\texorpdfstring{\textbf{3. Solution
Development}}{3. Solution Development}}\label{solution-development}

\textbf{Desktop Automation Tool:}\\
- Built with \textbf{PySimpleGUI} (user interface), \textbf{Beautiful
Soup/WebDriver} (web scraping), \textbf{PyAutoGUI} (GUI automation).\\
- \textbf{Workflow:}\\
1. Auto-login to portals/systems.\\
2. Extract transaction details and customer history.\\
3. Save consolidated data to a \textbf{Word document} (for analysts) and
\textbf{SQL Server} (for ML training).

\textbf{Significance of 50--60K Records:}\\
- A large dataset enables robust ML model training, capturing patterns
in both fraudulent and legitimate transactions.

\begin{center}\rule{0.5\linewidth}{0.5pt}\end{center}

\subsubsection{\texorpdfstring{\textbf{4. Machine Learning
Model}}{4. Machine Learning Model}}\label{machine-learning-model}

\paragraph{\texorpdfstring{\textbf{Data
Preparation}}{Data Preparation}}\label{data-preparation}

\begin{itemize}
\tightlist
\item
  \textbf{Label Conversion:} Convert labels (e.g., ``fraud''/``not
  fraud'') to numerical IDs (0/1).\\
\item
  \textbf{Train/Test Split:} 80\% training, 20\% validation/test.\\
\item
  \textbf{Dataset Format:} Use HuggingFace \texttt{Dataset} objects for
  seamless integration with ML libraries.
\end{itemize}

\paragraph{\texorpdfstring{\textbf{Tokenization}}{Tokenization}}\label{tokenization}

\begin{itemize}
\tightlist
\item
  \textbf{RoBERTa Tokenizer:} Breaks text (e.g., transaction notes) into
  subwords.\\
\item
  \textbf{Settings:}

  \begin{itemize}
  \tightlist
  \item
    \texttt{max\_length=512} (truncate/pad sequences to this length).\\
  \item
    \texttt{padding="max\_length"}, \texttt{truncation=True}
    (standardize input size).
  \end{itemize}
\end{itemize}

\paragraph{\texorpdfstring{\textbf{Model
Architecture}}{Model Architecture}}\label{model-architecture}

\begin{itemize}
\tightlist
\item
  \textbf{Base Model:} Pre-trained \textbf{RoBERTa} (exceles at text
  understanding).\\
\item
  \textbf{Classification Head:} Added on top to predict fraud
  probability (2 output classes).
\end{itemize}

\paragraph{\texorpdfstring{\textbf{Training
Configuration}}{Training Configuration}}\label{training-configuration}

\begin{itemize}
\tightlist
\item
  \textbf{Batch Size:} 16 (fits GPU memory).\\
\item
  \textbf{Learning Rate:} 2e-5 (avoids overwriting pre-trained
  knowledge).\\
\item
  \textbf{Epochs:} 3--4 (prevents overfitting).\\
\item
  \textbf{Mixed Precision Training:} Speeds up training with NVIDIA
  GPUs.
\end{itemize}

\paragraph{\texorpdfstring{\textbf{Metrics}}{Metrics}}\label{metrics}

\begin{itemize}
\tightlist
\item
  \textbf{Accuracy:} Overall correctness.\\
\item
  \textbf{Macro F1-Score:} Balances precision/recall across imbalanced
  classes.
\end{itemize}

\paragraph{\texorpdfstring{\textbf{Training
Process}}{Training Process}}\label{training-process}

\begin{itemize}
\tightlist
\item
  \textbf{HuggingFace Trainer API:} Handles gradient accumulation
  (simulates larger batches) and logging (track loss/metrics).
\end{itemize}

\paragraph{\texorpdfstring{\textbf{Model
Saving}}{Model Saving}}\label{model-saving}

\begin{itemize}
\tightlist
\item
  Save model weights, tokenizer, and label mappings for reproducibility.
\end{itemize}

\begin{center}\rule{0.5\linewidth}{0.5pt}\end{center}

\subsubsection{\texorpdfstring{\textbf{5. Validation
Process}}{5. Validation Process}}\label{validation-process}

\begin{enumerate}
\def\labelenumi{\arabic{enumi}.}
\tightlist
\item
  \textbf{Load Model:} Use saved weights to initialize the model.\\
\item
  \textbf{Test Data Preparation:} Tokenize test data identically to
  training.\\
\item
  \textbf{Custom Trainer:} Captures raw predictions (logits) for metric
  calculation.\\
\item
  \textbf{Prediction \& Metrics:}

  \begin{itemize}
  \tightlist
  \item
    \textbf{Classification Report:} Precision, recall, F1 per class.\\
  \item
    \textbf{Confusion Matrix:} Visualize false positives/negatives.
  \end{itemize}
\end{enumerate}

\begin{center}\rule{0.5\linewidth}{0.5pt}\end{center}

\subsubsection{\texorpdfstring{\textbf{6. Deployment and
Monitoring}}{6. Deployment and Monitoring}}\label{deployment-and-monitoring}

\textbf{Deployment Steps:}\\
1. Package model into a Docker container.\\
2. Expose as an API (e.g., FastAPI/Flask) for real-time predictions.\\
3. Deploy on cloud (AWS/GCP) with auto-scaling.

\textbf{Monitoring:}\\
- Track \textbf{latency}, \textbf{throughput}, and \textbf{model drift}
(e.g., Prometheus/Grafana).\\
- Log predictions to detect anomalies (e.g., sudden drop in F1-score).

\begin{center}\rule{0.5\linewidth}{0.5pt}\end{center}

\subsubsection{\texorpdfstring{\textbf{7. Handling Data Drift and
Re-Training}}{7. Handling Data Drift and Re-Training}}\label{handling-data-drift-and-re-training}

\begin{itemize}
\tightlist
\item
  \textbf{Detect Drift:} Statistical tests (e.g., Kolmogorov-Smirnov) on
  feature distributions.\\
\item
  \textbf{Re-Training:}

  \begin{enumerate}
  \def\labelenumi{\arabic{enumi}.}
  \tightlist
  \item
    Triggered when drift exceeds a threshold.\\
  \item
    Fine-tune model on new data + historical data to retain prior
    knowledge.
  \end{enumerate}
\end{itemize}

\begin{center}\rule{0.5\linewidth}{0.5pt}\end{center}

\subsubsection{\texorpdfstring{\textbf{8. Scaling the
Solution}}{8. Scaling the Solution}}\label{scaling-the-solution}

\begin{itemize}
\tightlist
\item
  \textbf{Horizontal Scaling:} Add more API servers behind a load
  balancer.\\
\item
  \textbf{Batch Processing:} Use Spark/Dask for large overnight
  inference jobs.\\
\item
  \textbf{Database Optimization:} Sharding/partitioning in SQL Server
  for faster queries.
\end{itemize}

\begin{center}\rule{0.5\linewidth}{0.5pt}\end{center}

\subsubsection{\texorpdfstring{\textbf{Example
Workflow}}{Example Workflow}}\label{example-workflow}

\begin{enumerate}
\def\labelenumi{\arabic{enumi}.}
\tightlist
\item
  An analyst opens the desktop tool, which auto-fetches transaction
  data.\\
\item
  The ML model assigns a risk score (e.g., 0.95 = high risk).\\
\item
  High-risk cases are prioritized; low-risk cases are auto-approved.\\
\item
  Analysts focus on 10--15 high-risk alerts/day instead of 40 low-value
  ones.
\end{enumerate}

\subsubsection{\texorpdfstring{\textbf{Conclusion}}{Conclusion}}\label{conclusion}

By automating data collection and deploying an ML model, the solution
reduces review time by \textbf{60--70\%}, clears backlogs, and adapts to
evolving fraud patterns through continuous monitoring. This end-to-end
pipeline transforms a reactive process into a proactive, scalable
system.

\end{document}
